%!TEX root = ../main.tex
\newpage
\section{見かけの位置へのヤコビアン$J_app$の導出}\label{sec:appendix-jacobian}

\cref{sec:position-correction}で述べた位置補正において、共分散補正(\cref{sec:covariance-correction})等で用いる変換の勾配には、真の位置$\bm{p}(r,z)$から見かけの位置$\bm{p}'(r',z')$へのヤコビアン行列$J_app$が必要となる。
本節では、$rz$平面における$2 \times 2$のヤコビアン
\begin{equation}\label{eq:Jacobian-def-rz}
  J_{rz} = \frac{\partial (r', z')}{\partial (r, z)} = \begin{pmatrix} \dfrac{\partial r'}{\partial r} & \dfrac{\partial r'}{\partial z} \\[1em] \dfrac{\partial z'}{\partial r} & \dfrac{\partial z'}{\partial z} \end{pmatrix}
\end{equation}
の各成分を、\cref{sec:appendix-apparent-depth}で得られた見かけの位置の式\eqref{eq:apparent-position}に基づき導出する。
$r'$、 $z'$は$\theta_a$、 $\theta_w$の関数として表されており、これらは$(r,z)$に依存するため、連鎖律を用いて角度の偏微分を経由して求める。

\subsection{角度の偏微分係数}

幾何学的拘束およびスネルの法則から、$\theta_w$の$r$, $z$に関する偏微分係数を求める。
カメラと水面交点$(s, 0)$および\cref{eq:-r-s}より
\begin{equation}\label{eq:geom-constraint}
  r = s - z \tan \theta_w, \qquad s = H \tan \theta_a
\end{equation}
であり、$s$を消去すると
\begin{equation}\label{eq:constraint-rz}
  r = H \tan \theta_a - z \tan \theta_w
\end{equation}
が成り立つ。
この両辺の全微分をとると
\begin{equation}\label{eq:total-diff-rz}
  dr = H \sec^2 \theta_a \, d\theta_a - \tan \theta_w \, dz - z \sec^2 \theta_w \, d\theta_w
\end{equation}
となる。
一方、スネルの法則\cref{eq:snell}を微分すると
\begin{equation}\label{eq:snell-diff}
  \cos \theta_a \, d\theta_a = n \cos \theta_w \, d\theta_w \quad \Rightarrow \quad d\theta_a = n \, \frac{\cos \theta_w}{\cos \theta_a} \, d\theta_w
\end{equation}
である。
\eqref{eq:total-diff-rz}に\eqref{eq:snell-diff}を代入し、$d\theta_w$について整理すると
\begin{equation}\label{eq:dtheta-w-expr}
  dr + \tan \theta_w \, dz = \left( H \sec^2 \theta_a \cdot n \, \frac{\cos \theta_w}{\cos \theta_a} - z \sec^2 \theta_w \right) d\theta_w
\end{equation}
となる。
右辺の$d\theta_w$の係数を
\begin{equation}\label{eq:Delta-def}
  \Delta \equiv n H \, \frac{\cos \theta_w}{\cos^3 \theta_a} - z \sec^2 \theta_w
\end{equation}
と定義する。
$z < 0$であることから第2項は正となり、$\Delta \neq 0$が保たれる。
これより、\eqref{eq:dtheta-w-expr}を$d\theta_w$について解く（両辺を$\Delta$で割る）と
\begin{equation}\label{eq:dtheta-w-solved}
  d\theta_w = \frac{1}{\Delta}dr + \frac{\tan\theta_w}{\Delta}dz
\end{equation}
を得る。
一方で、$\theta_w=\theta_w(r,z)$の全微分は一般に
\begin{equation}
  d\theta_w = \left(\frac{\partial \theta_w}{\partial r}\right)dr + \left(\frac{\partial \theta_w}{\partial z}\right)dz
\end{equation}
と書けるため、$dr$および$dz$の係数を比較すれば、$\theta_w$の偏微分係数は
\begin{equation}\label{eq:partial-theta-w}
  \frac{\partial \theta_w}{\partial r} = \frac{1}{\Delta}, \qquad \frac{\partial \theta_w}{\partial z} = \frac{\tan \theta_w}{\Delta}
\end{equation}
と得られる。
また、以降の記述を簡潔にするため
\begin{equation}\label{eq:alpha-def}
  \beta \equiv n \, \frac{\cos \theta_w}{\cos \theta_a}
\end{equation}
とおく。

\subsection{ヤコビアン各成分の導出}

見かけの位置は\eqref{eq:r'}, \eqref{eq:z'}で与えられる:
\begin{equation}\label{eq:r'-z'-recap}
  r' = r + (n^2 - 1) z \tan^3 \theta_w, \qquad z' = \frac{z}{n} \left( \frac{\cos \theta_a}{\cos \theta_w} \right)^3
\end{equation}

\subsubsection{\boldmath{$J_{11} = \dfrac{\partial r'}{\partial r}$}}

\eqref{eq:r'-z'-recap}の$r'$を$r$で偏微分する。
$r$は直接現れるほか、$\theta_w(r,z)$を通じて間接的に現れるため、連鎖律より:
\begin{equation}\label{eq:dr'-dr}
  \frac{\partial r'}{\partial r} = 1 + (n^2 - 1) z \cdot \frac{\partial}{\partial \theta_w}(\tan^3 \theta_w) \cdot \frac{\partial \theta_w}{\partial r} = 1 + (n^2 - 1) z \cdot 3 \tan^2 \theta_w \sec^2 \theta_w \cdot \frac{1}{\Delta}
\end{equation}

したがって:
\begin{equation}\label{eq:J11}
  \frac{\partial r'}{\partial r} = 1 + \frac{3(n^2 - 1) z \tan^2 \theta_w \sec^2 \theta_w}{\Delta}
\end{equation}


\subsubsection{\boldmath{$J_{12} = \dfrac{\partial r'}{\partial z}$}}

$r'$を$z$で偏微分する。
積の微分と$\dfrac{\partial \theta_w}{\partial z} = \dfrac{\tan \theta_w}{\Delta}$を用いると:
\begin{align}
  \frac{\partial r'}{\partial z} &= (n^2 - 1) \left[ \tan^3 \theta_w + z \cdot 3 \tan^2 \theta_w \sec^2 \theta_w \cdot \frac{\tan \theta_w}{\Delta} \right] \notag \\
  &= (n^2 - 1) \tan^3 \theta_w \left( 1 + \frac{3 z \sec^2 \theta_w}{\Delta} \right)
\end{align}
したがって:
\begin{equation}\label{eq:J12}
  \frac{\partial r'}{\partial z} = (n^2 - 1) \tan^3 \theta_w \left( 1 + \frac{3 z \sec^2 \theta_w}{\Delta} \right)
\end{equation}


\subsubsection{\boldmath{$J_{21} = \dfrac{\partial z'}{\partial r}$}}

$z' = \dfrac{z}{n} \left( \dfrac{\cos \theta_a}{\cos \theta_w} \right)^3$に対し、対数微分を用いる。
\begin{equation}
  \ln z' = \ln z - \ln n + 3(\ln \cos \theta_a - \ln \cos \theta_w)
\end{equation}
を$r$で偏微分すると、$\ln z$の項は$r$に依存しないため:
\begin{equation}
  \frac{1}{z'} \frac{\partial z'}{\partial r} = 3 \left( -\tan \theta_a \frac{\partial \theta_a}{\partial r} + \tan \theta_w \frac{\partial \theta_w}{\partial r} \right)
\end{equation}

$\dfrac{\partial \theta_a}{\partial r} = \beta \dfrac{\partial \theta_w}{\partial r}$および\eqref{eq:partial-theta-w}より:
\begin{equation}
  \frac{\partial z'}{\partial r} = 3 z' \left( \tan \theta_w - \beta \tan \theta_a \right) \frac{\partial \theta_w}{\partial r} = \frac{3 z'}{\Delta} \left( \tan \theta_w - \beta \tan \theta_a \right)
\end{equation}
ここで$\beta = n \cos \theta_w / \cos \theta_a$および$\sin \theta_a = n \sin \theta_w$を用いると:
\begin{equation}
  \tan \theta_w - \beta \tan \theta_a = \tan \theta_w \left( 1 - n^2 \, \frac{\cos^2 \theta_w}{\cos^2 \theta_a} \right)
\end{equation}
と整理できる。
よって:
\begin{equation}\label{eq:J21}
  \frac{\partial z'}{\partial r} = \frac{3 z' \tan \theta_w}{\Delta} \left( 1 - n^2 \, \frac{\cos^2 \theta_w}{\cos^2 \theta_a} \right)
\end{equation}

\subsubsection{\boldmath{$J_{22} = \dfrac{\partial z'}{\partial z}$}}

同様に対数微分した式を$z$で偏微分する。
$\ln z$の項から$\dfrac{1}{z}$が現れる:
\begin{equation}
  \frac{1}{z'} \frac{\partial z'}{\partial z} = \frac{1}{z} + 3 \left( \tan \theta_w - \beta \tan \theta_a \right) \frac{\partial \theta_w}{\partial z} = \frac{1}{z} + 3 \left( \tan \theta_w - \beta \tan \theta_a \right) \frac{\tan \theta_w}{\Delta}
\end{equation}
$J_{21}$の導出と同様に括弧内を整理し:
\begin{equation}\label{eq:J22}
  \frac{\partial z'}{\partial z} = \frac{z'}{z} + \frac{3 z' \tan^2 \theta_w}{\Delta} \left( 1 - n^2 \, \frac{\cos^2 \theta_w}{\cos^2 \theta_a} \right)
\end{equation}

\vspace{1em}
以上より、求めるヤコビアン$J_{\mathrm{app}}$（$rz$平面における変換$\bm{p}=(r,z)^\top \mapsto \bm{p}'=(r',z')^\top$の勾配）は
\begin{equation}\label{eq:Jacobian-app-matrix}
  J_{\mathrm{app}} = \frac{\partial (r', z')}{\partial (r, z)}
  =
  \begin{pmatrix}
    \dfrac{\partial r'}{\partial r} & \dfrac{\partial r'}{\partial z} \\[1.2em]
    \dfrac{\partial z'}{\partial r} & \dfrac{\partial z'}{\partial z}
  \end{pmatrix}
  =
  \begin{pmatrix}
    \displaystyle 1 + \frac{3(n^2 - 1) z \tan^2 \theta_w \sec^2 \theta_w}{\Delta}
    &
    \displaystyle (n^2 - 1) \tan^3 \theta_w \left( 1 + \frac{3 z \sec^2 \theta_w}{\Delta} \right)
    \\[1.5em]
    \displaystyle \frac{3 z' \tan \theta_w}{\Delta} \left( 1 - n^2 \, \frac{\cos^2 \theta_w}{\cos^2 \theta_a} \right)
    &
    \displaystyle \frac{z'}{z} + \frac{3 z' \tan^2 \theta_w}{\Delta} \left( 1 - n^2 \, \frac{\cos^2 \theta_w}{\cos^2 \theta_a} \right)
  \end{pmatrix}
\end{equation}
で与えられる。

\subsection{数値実装上の注意：中心軸および水面近傍の極限形}
\label{sec:appendix-jacobian-limit}

数値計算の実装においては、$\Delta$が小さい場合や$\tan\theta_w$が極めて小さい場合に、分母の極小化や桁落ちにより不安定化することがある。
そのため、中心軸付近（$r\to 0$）および水面付近（$z\to 0$）では、適当な閾値$\epsilon$を設けて極限形に切り替える処理が有効である。

\paragraph{\boldmath{$r\to 0$（ナディア視点）}}
$r\to 0$では$\theta_w\to 0$かつ$\theta_a\to 0$となり、
$\tan\theta_w\to 0$、$\cos\theta_w\to 1$、$\cos\theta_a\to 1$より$\Delta\to nH-z$、さらに$z'\to z/n$である。
したがって、行列要素のうち$\tan\theta_w$を含む項は消失し、
\begin{equation}\label{eq:Jacobian-app-r0}
  \lim_{r\to 0} J_{\mathrm{app}}
  =
  \begin{pmatrix}
    1 & 0 \\
    0 & \dfrac{1}{n}
  \end{pmatrix}
\end{equation}
を得る。

\paragraph{\boldmath{$z\to 0$（水面付近）}}
$z\to 0$では$z'\to 0$である一方、角度$\theta_w,\theta_a$は（$z=0$における幾何学拘束$r=H\tan\theta_a$およびスネルの法則）から決まる有限値に収束する。
このとき$J_{11}$の補正項は$z$に比例して消失するため$J_{11}\to 1$、また$J_{21}$および$J_{22}$の第2項は$z'$に比例して消失するため$J_{21}\to 0$となる。
よって
\begin{equation}\label{eq:Jacobian-app-z0}
  \lim_{z\to 0} J_{\mathrm{app}}
  =
  \begin{pmatrix}
    1 & (n^2-1)\tan^3\theta_w \\
    0 & \dfrac{1}{n}\left(\dfrac{\cos\theta_a}{\cos\theta_w}\right)^3
  \end{pmatrix}
\end{equation}
である。

